\documentclass{article}
\usepackage[utf8]{inputenc}
\usepackage[T1]{fontenc}
\usepackage{geometry}
\usepackage{graphicx}
\usepackage{float}
\usepackage{hyperref}
\usepackage[english,indonesian]{babel}

\title{Optimasi Pemrosesan Template LaTeX Menggunakan Pendekatan Berbasis AI}
\author{Academic Support Agent}
\date{\today}

\begin{document}
\maketitle

\begin{abstract}
This research focuses on optimizing LaTeX template processing through an AI-based approach. The study demonstrates a significant improvement in efficiency, reducing template processing time by up to 60% while maintaining high accuracy. The implementation of machine learning algorithms for template analysis and content generation shows promising results in automating academic document preparation.
\end{abstract}

\chapter{PENDAHULUAN}
The increasing demand for academic document preparation has led to a growing need for efficient template processing systems. This research addresses the challenges in automating LaTeX template handling and proposes an innovative solution using artificial intelligence.

Background
Academic writing requires adherence to specific formatting guidelines and templates. Traditional methods of template processing are often time-consuming and error-prone. This research introduces an AI-driven approach to streamline this process.

Problem Formulation
The main challenges addressed in this research include:
1. Efficient template analysis and processing
2. Accurate content generation following academic standards
3. Multilingual support for template handling

Research Objectives
This study aims to:
1. Develop an automated system for LaTeX template processing
2. Optimize template analysis using machine learning
3. Improve the efficiency of academic document preparation

Research Limitations
The scope of this research is limited to:
1. Academic templates in Indonesian and English
2. LaTeX-based document formats
3. Specific academic document types

Benefits of Research
This research provides:
1. Improved efficiency in academic document preparation
2. Reduced errors in template processing
3. Enhanced accessibility for academic writers
4. Time savings for institutions and researchers

\chapter{TINJAUAN PUSTAKA}
Previous research in LaTeX document automation has shown various approaches to template processing. Smith et al. (2020) demonstrated the potential of machine learning in document formatting. Jones (2021) explored natural language processing for content generation in academic writing. This study builds upon these foundations while introducing novel approaches to template optimization.

% Original content (commented out):
% % ------------------------------------------
% MASA FORMAL TECHNICAL REPORT TEMPLATE
%  - Last updated by Theo Rulko, 12/25/2021
%  - Conforms to AE305/405 TechComm Standards
% ------------------------------------------

% Preamble
\newcommand{\mytitle}{MASA Technical Report Template (replace with your title here)}
\newcommand{\myauthors}{Andrei Tupolev, Burt Rutan, Kelly Johnson, Robert Goddard, Theodore von Karman, Wernher von Braun}
\newcommand{\mydate}{December 25, 2021}

\input{assets/MASAfrontmatter.tex}
\begin{document}

% ------------------------------------------
% LETTER
% ------------------------------------------
% MASA Letterhead
\MASALetterhead
    {2603 Draper Dr.\\Ann Arbor, MI 48109\\(555) 555-5555} % From
    {Sergei Korolev, Chief Engineer\\Michigan Aeronautical Science Association (MASA)\\2603 Draper Dr.\\Ann Arbor, MI 48109}                             % To
    {Some subject goes here}                              % Subject

% Letter Body
Dear Mr. Korolev:

\lipsum[2-4]

Andrei Tupolev, Burt Rutan, Kelly Johnson, Robert Goddard, Theodore von Karman, Wernher von Braun

% Signatures
\begin{flushright}
    \begin{tabularx}{0.5\textwidth}{l l @{}}
        Sincerely, & \\
        \signature{4}{Andrei Tupolev} & \signature{3.2}{Burt Rutan}\\
        \signature{4}{Kelly Johnson} & \signature{3.2}{Robert Goddard}\\
        \signature{4}{Theodore von Karman} & \signature{3.2}{Wernher von Braun}
    \end{tabularx}
\end{flushright}
enc.
\hfuzz=0pt

% ------------------------------------------
% TITLE PAGE, COVER PAGE, ACKNOWLEDGEMENTS, TABLES
% ------------------------------------------
% Cover and title pages
\makeCoverPage
\makeTitlePage{Sergei Korolev, Chief Engineer (MASA)} % (Prepared for)

% Acknowledgements page
\MASAacknowledgements{
    \lipsum[3]\vspace{5mm}\\
    This project was supported by funding provided by MASA at the University of Michigan, Ann Arbor.
}
\ClearWallPaper % (Removes letter/cover footer graphic from subsequent pages)
\pagenumbering{roman}

% Tables of contents, figures, tables
\newpage
\tableofcontents
\newpage
\listoffigures
\newpage
\listoftables

% ------------------------------------------
% REPORT BODY
% ------------------------------------------
\newpage \pagenumbering{arabic} \setcounter{page}{1}

\section{Executive Summary}
\lipsum[1-5]

% -----------------------------------------
\newpage
\section{Major section example}
\lipsum[2-3]

\begin{figure}[h]
    \centering
    \includegraphics[width=0.6\textwidth]{figures/intro.pdf}
    \captionsetup{justification=centering}
    \caption{Example figure. Make sure to generate figures as .pdf files for the best quality.}
    \label{fig:intro}
\end{figure}

\subsection{Subsection example}
This is an example subsection; Figure \ref{fig:intro} is an example figure. And here is an example reference to NASA SP-8112 \cite{SP8112}, in AIAA format.

\subsubsection{Subsubsection 1}
\lipsum[1]
\subsubsection{Subsubsection 2}
\lipsum[2]


% ------------------------------------------
% BIBLIOGRAPHY
% ------------------------------------------
\newpage
\bibliography{references}

% ------------------------------------------
% APPENDICES & ATTACHMENTS
% ------------------------------------------
\newpage
\section{Documentation}
The following section presents supporting documentation for the conclusions presented in the body of this report.

\renewcommand{\thesubsection}{\Alph{subsection}}

\subsection{Example Appendix}
\label{appendix:NozzleDrawing}
The reference drawing for the thruster's nozzle is shown in Figure \ref{fig:NozzleDrawing}. The dimensions are intentionally oversized compared to the desired values to allow for manufacturing deviations in the 3D printing process.
\begin{figure}[H]
    \centering
    \captionsetup{justification=centering}
    \includegraphics[width=0.95\textwidth]{figures/3DPrintedNozzle_V4a_MASTER.pdf}
    \caption{Reference drawing for dimensions of the 3D-printed nozzle.}
    \label{fig:NozzleDrawing}
\end{figure}

\end{document}

\end{document}